\chapter*{Glossaire}
\addcontentsline{toc}{chapter}{Glossaire}

Dans ce section , nous allons définir tous les termes techniques et acronymes que nous avons utilisés dans notre rapport.

\begin{itemize}
    \item \textbf{Régression linéaire} : une technique d'analyse de données qui prédit la valeur de données inconnues en utilisant une autre valeur de données apparentée et connue.
    \item \textbf{Classification} : une technique d'analyse de données qui prédit la catégorie ou la classe à laquelle appartient une observation donnée.
    \item \textbf{Régression logistique} : une technique de classification qui utilise une fonction logistique pour modéliser la probabilité qu'une observation appartienne à une classe particulière.
    \item \textbf{Decision tree} : un algorithme d'apprentissage automatique qui utilise une multitude d'arbres de décision pour améliorer la précision de la classification ou de la régression.
    \item \textbf{Arbre de décision} : un modèle prédictif qui utilise une structure arborescente pour représenter les décisions et leurs conséquences.
    \item \textbf{SVC} : Support Vector Classifier, un algorithme de classification qui trouve l'hyperplan optimal pour séparer les différentes classes dans un espace de caractéristiques.
    \item \textbf{hyperplan} : une frontière de décision qui sépare les différentes classes dans un espace de caractéristiques.
    \item \textbf{KNN} : K Nearest Neighbors, un algorithme de classification qui classe une observation en fonction des classes de ses K voisins les plus proches dans l'espace de caractéristiques.
    \item \textbf{univariée} : une analyse qui examine une seule variable à la fois.
    \item \textbf{bivariée} : une analyse qui examine la relation entre deux variables.
    \item \textbf{valeurs manquantes} : des données absentes dans un ensemble de données.
    \item \textbf{bibliothèque} : un ensemble de fonctions et de routines pré-écrites qui peuvent être utilisées pour accomplir des tâches spécifiques dans un langage de programmation.
    \item \textbf{pandas} : une bibliothèque Python utilisée pour la manipulation et l'analyse de données.
    \item \textbf{matplotlib} : une bibliothèque Python utilisée pour la création de visualisations statiques, animées et interactives.
    \item \textbf{seaborn} : une bibliothèque Python basée sur matplotlib qui fournit une interface de haut niveau pour dessiner des graphiques statistiques attrayants et informatifs.
    \item \textbf{numpy} : une bibliothèque Python utilisée pour le calcul scientifique et  le traitement de tableaux multidimensionnels.
    \item \textbf{score validation} : une mesure de la performance d'un modèle sur un ensemble de données de test.
    \item \textbf{score entraînement} : une mesure de la performance d'un modèle sur l'ensemble de données d'entraînement.
    \item \textbf{variable quantitative} : une variable qui peut être mesurée numériquement.
    \item \textbf{variable discrète} : une variable qui peut prendre un nombre fini ou dénombrable de valeurs.
    \item \textbf{variable cible} : la variable que l'on cherche à prédire ou à expliquer dans une analyse de données.
    \item \textbf{dataset} : un ensemble structuré de données, souvent présenté sous forme de tableau.
    \item \textbf{scraper} : le processus d'extraction de données à partir de sites web.
    \item \textbf{Kaggle} : une plateforme en ligne pour les compétitions de science des données et le partage de datasets.
    \item \textbf{feature} : une caractéristique ou une variable utilisée comme entrée dans un modèle d'apprentissage automatique.
    \item \textbf{scikit-learn} : une bibliothèque Python pour l'apprentissage automatique.
    \item \textbf{StandardScaler} : une technique de normalisation des données qui transforme les caractéristiques pour qu'elles aient une moyenne de 0 et un écart type de 1.
    \item \textbf{train-test split} : la division d'un ensemble de données en un ensemble d'entraînement et un ensemble de test pour évaluer les performances d'un modèle.
    \item \textbf{accuracy} : la proportion de prédictions correctes parmi le nombre total de prédictions.
    \item \textbf{precision} : la proportion de vrais positifs parmi les prédictions positives.
    \item \textbf{recall} : la proportion de vrais positifs parmi les observations réelles positives.
    \item \textbf{F1-score} : la moyenne harmonique entre la précision et le recall.
    \item \textbf{courbe d'apprentissage} : un graphique qui montre comment les performances d'un modèle évoluent avec la quantité de données d'entraînement.
    \item \textbf{modèle} : une représentation mathématique ou algorithmique utilisée pour faire des prédictions ou des classifications basées sur des données.
    \item \textbf{surapprentissage} : un phénomène où un modèle apprend trop bien les détails et le bruit des données d'entraînement, ce qui nuit à sa capacité à généraliser à de nouvelles données.
\end{itemize}